% Generated draft: requires mathematical and source review.
\documentclass[11pt]{article}
\usepackage[margin=25mm]{geometry}
\usepackage{fontspec}
\setmainfont{Latin Modern Roman}
\usepackage{amsmath,amssymb,mathtools,unicode-math}
\setmathfont{Latin Modern Math}
\usepackage{enumitem,needspace,xurl,hyperref}
\setlength{\parindent}{0pt}
\setlength{\parskip}{5pt}
\setlength{\emergencystretch}{3em}
\newcommand{\N}{\mathbb{N}}
\newcommand{\Z}{\mathbb{Z}}
\newcommand{\Q}{\mathbb{Q}}
\newcommand{\R}{\mathbb{R}}
\newcommand{\C}{\mathbb{C}}
\newcommand{\F}{\mathbb{F}}
\newcommand{\A}{\mathbb{A}}
\newcommand{\PP}{\mathbb{P}}
\newcommand{\E}{\mathbb{E}}
\newcommand{\eps}{\varepsilon}
\newcommand{\dd}{\,\mathrm{d}}
\begin{document}
\section*{UnsolvedMath ID: 701142}
\textbf{Source problem code:} catalogue entry
\subsection{Definitions and mathematical statement}
A group is periodic if every element has finite order. An involution is an element of order exactly two. A group is locally cyclic if each of its finitely generated subgroups is cyclic, and a \(2\)-group here means a group in which every element has order a power of \(2\), including \(2^0=1\).

Let \(G\) be any periodic group, and let \(i\in G\) be an involution such that its centralizer
\[
C_G(i)=\{g\in G:gi=ig\}
\]
is a locally cyclic \(2\)-group. Define
\[
S_i=\{g\in G:\operatorname{ord}(g)\text{ is odd and }i^{-1}gi=g^{-1}\}.
\]
The identity belongs to \(S_i\), since its order is \(1\). Must \(S_i\) be a subgroup of \(G\) for every such pair \((G,i)\)? No finiteness assumption on \(G\) is imposed beyond periodicity.
\subsection{Short English statement}
In a periodic group, suppose an involution has a locally cyclic centralizer consisting of elements of \(2\)-power order. Must all odd-order elements that this involution inverts, including the identity, together form a subgroup?
\subsection{Sources}
\begin{itemize}
\item Local UnsolvedMath registry, permanent ID \(701142\), exact source code \texttt{catalogue entry}, title \emph{Kourovka Notebook Problem 15.54}; supplied source record, upstream revision \texttt{d371ee917ac9647fbb06cd6fc529deec6f3dda82}.
\item Attributed in the supplied record to V. D. Mazurov, proposed in 2002: \emph{The Kourovka Notebook (Unsolved Problems in Group Theory)}, 15th issue (2002), Problem 15.54. The record also identifies the 21st issue edition (2026), arXiv:1401.0300v48. The supplied collection description names E. I. Khukhro and V. D. Mazurov as editors.
\item The registry labels the problem \texttt{open}. This is a supplied source label, not an independently verified assessment of the mathematical literature. The record was updated on \texttt{2026-10-16}; this timestamp is registry metadata, not a supplied literature-verification date. No source was independently consulted for this draft.
\end{itemize}
\end{document}
